% 第九卷 14.3.3：三种计划与行动组织方式（跨栏，约160 mm宽）。
% FigureModel：语言模型调用；FigureOutput：工具执行与调度；虚线：变量引用或重新规划。
\begin{tikzpicture}[
  x=1mm,y=1mm,
  figure picture,
  model/.style={figure node,draw=FigureModel,fill=FigureModelFill,
    minimum width=38mm,align=center,minimum height=11mm,inner xsep=2pt,line width=.8pt},
  tool/.style={figure node,draw=FigureOutput,fill=FigureOutputFill,
    minimum height=8mm,inner xsep=2pt},
  title/.style={font=\figurefont\small,text=FigureInk,align=center},
  ref/.style={figure flow,dash pattern=on 2pt off 1.5pt}
]
  % (a) ReAct
  \node[title] at (22,64) {（a）ReAct：逐步交替};
  \node[model] (ra) at (22,48) {语言模型\\Thought$_t$与Action$_t$};
  \node[tool,minimum width=30mm] (rt) at (22,22) {工具执行};
  \draw[figure flow] ([xshift=-7mm]ra.south) --
    node[figure edge label,left=1pt] {Action$_t$} ([xshift=-7mm]rt.north);
  \draw[figure flow] ([xshift=7mm]rt.north) --
    node[figure edge label,right=1pt] {Observation$_t$} ([xshift=7mm]ra.south);
  \node[figure note] at (22,7) {$T$次工具调用需$T{+}1$次模型调用\\每次重读提示与全部历史};

  % (b) ReWOO
  \node[title] at (74,64) {（b）ReWOO：先规划后执行};
  \node[model] (wp) at (74,48) {规划器\\一次写出含\#E$_k$的计划};
  \node[tool,minimum width=11mm] (w1) at (60,31) {\#E$_1$};
  \node[tool,minimum width=11mm] (w2) at (74,31) {\#E$_2$};
  \node[tool,minimum width=11mm] (w3) at (88,31) {\#E$_3$};
  \node[model] (ws) at (74,14) {求解器\\读计划与证据作答};
  \foreach \w in {w1,w2,w3} {
    \draw[figure flow] (wp.south -| \w) -- (\w.north);
    \draw[figure flow] (\w.south) -- (ws.north -| \w);
  }
  \draw[ref] (w1.east) -- (w2.west);
  \node[figure note] at (74,3) {工作器按变量引用代入结果\\共2次模型调用};

  % (c) LLMCompiler
  \node[title] at (126,64) {（c）LLMCompiler：依赖图并行};
  \node[model] (cp) at (126,48) {函数调用规划器\\输出任务依赖图};
  \node[tool,minimum width=34mm,align=center] (cf) at (126,34) {任务获取单元\\依赖满足即派发};
  \node[tool,minimum width=15mm] (c1) at (117,21) {执行器};
  \node[tool,minimum width=15mm] (c2) at (135,21) {执行器};
  \node[model] (cj) at (126,7) {汇合\\作答或重新规划};
  \draw[figure flow] (cp.south) -- (cf.north);
  \draw[figure flow] (cf.south -| c1) -- (c1.north);
  \draw[figure flow] (cf.south -| c2) -- (c2.north);
  \draw[figure flow] (c1.south) -- (cj.north -| c1);
  \draw[figure flow] (c2.south) -- (cj.north -| c2);
  \draw[ref] (cj.east) -- (148,7) -- node[figure edge label,right=1pt,text width=8mm] {重新\\规划} (148,48) -- (cp.east);
\end{tikzpicture}
